\section{Validazione del corretto funzionamento dell’ambiente}

La validazione iniziale è stata impostata come \emph{smoke test} progressivo: prima esempi base Plexe/Veins, poi scenari con eterogeneità di interfacce, infine verifica dei sotto-progetti.

I tutorial ufficiali forniscono configurazioni utili per un controllo rapido della pipeline end-to-end \cite{plexe_tutorial,plexe_example6}. I comandi utilizzati sono del tipo:

\begin{lstlisting}[caption={Esecuzione di scenari tutorial da CLI (Cmdenv)},label={lst:smoke_commands}]
cd /path/to/plexe/examples/platooning
./run -u Cmdenv -c BrakingScenario
./run -u Cmdenv -c BrakingScenarioWithApp

cd /path/to/plexe/examples/plexe_veins
./run -u Cmdenv -c LanradioScenario
./run -u Cmdenv -c HeterogeneousScenario
\end{lstlisting}

In questa fase la verifica non è solo ``la simulazione parte'', ma include almeno:
\begin{itemize}
\item assenza di errori di linkage o moduli OMNeT++ non risolti;
\item handshake corretto TraCI tra OMNeT++ e SUMO;
\item generazione dei file \texttt{.sca}/\texttt{.vec} e log coerenti con la configurazione selezionata;
\item comportamento dinamico plausibile del plotone (assenza di artefatti macroscopici dovuti a mismatch di configurazione).
\end{itemize}

Per ridurre il tempo di diagnosi, è stato utile introdurre un test automatico minimo che esegue un set ristretto di configurazioni e fallisce alla prima anomalia:

\begin{lstlisting}[caption={Script minimale di smoke validation},label={lst:smoke_script}]
#!/usr/bin/env bash
set -euo pipefail

SCENARIOS=("BrakingScenario" "BrakingScenarioWithApp")
LOGDIR="logs/smoke"
mkdir -p "$LOGDIR"

for cfg in "${SCENARIOS[@]}"; do
  echo "[SMOKE] running $cfg"
  ./run -u Cmdenv -c "$cfg" -r 0 > "$LOGDIR/$cfg.log" 2>&1
done

echo "[SMOKE] completed"
\end{lstlisting}
